\section{Exact integer multiplication and its cost}
\label{sec:assembly}

We now use the preceding algorithms to multiply two $n$-bit integers.
An equal-width address swap on logical bit volume $V$ costs
$O(Vu^\tau)$ for width $u$.  A completed parallel normalized butterfly
layer costs $O(Vd^{\lambda'})$ and has coefficient error less than
$\sqrt2\,2^{-p}$.  The ordered Gaussian resampling maps are
contractions; their disk-grid approximations have error less than
$p^2 2^{-p}$ and cost $O(tp^{3/2+\delta}\alpha)$ on a line of length
$t$.  Address permutations and intermediate dyadic arithmetic are
exact.  We choose compatible parameters for these procedures, recover
the integer coefficients from their approximations, and account for the
remaining tape operations.

\subsection{A fixed rational choice}

The two motif exponents were fixed in \eqref{eq:explicit-motif-exponents}.
Take the following fixed rational numbers:
\begin{equation}\label{eq:fixed-parameters}
\begin{gathered}
 \tau=\sigma=1-2^{-50},\qquad \beta=\frac12,\qquad \delta=\frac1{16},\\
 \lambda=1-2^{-52},\qquad c=2^{-56},\qquad
 \lambda'=1-2^{-54},\\
 C_1=20,\qquad \epsilon=2^{-75},\qquad \kappa=2^{-182}.
\end{gathered}
\end{equation}
These constants are deliberately conservative; no optimization is claimed.

The layer conditions hold with room to spare:
\[
 \max\{\tau,\sigma\}<\lambda<1,\qquad
 \tau(1+c/\beta)<1-31\cdot2^{-55}<\lambda,
\]
and
\[
 \sigma+\beta(1-\sigma)=1-2^{-51}<\lambda<\lambda'<1.
\]
The strict final inequality absorbs the logarithmic number of groups in
a base-$m$ partition.  The dimension exponent also satisfies
\[
\begin{gathered}
 \epsilon<\frac1{12},\quad \epsilon C_1<1,\quad
 \epsilon(2-\tau)<1-\tau,\quad
 \frac34+\delta+\frac32\epsilon<1,\\
 \epsilon(1+c)<1,\qquad \epsilon+\delta<1.
\end{gathered}
\]

\subsection{Input and transform sizes}

The input is the promised string $x\#y$, with both operands written
most significant bit first.  In one left-to-right pass, copy $x$ and
$y$ to two fixed work tapes and count the bits of $x$ with a binary
counter.  The promise gives both operands this counted length $n$.
The counter bound in Section~\ref{sec:streams} and the two sequential
copies give cost $O(n)$, including all leading zeroes.  The work-copy
heads finish just after the least significant bits, ready to move
leftward when the digits are formed.

\subsubsection{The padded box and working precision}

After obtaining $n$, compute the size parameters below with the
separate setup costs stated in this subsection.  For sufficiently
large $n$, put
\[
 b=\lceil\log_2n\rceil,\quad p=6b,\quad d=\lfloor b^\epsilon\rfloor,
 \quad K=\lfloor d^c\rfloor.
\]
Take the power of two $T$ in $[4n/b,8n/b)$, and let
$r=2^{\lceil(\log_2T)/d\rceil}$ and $\ell=\log_2r$.  Choose
$d\ell-\log_2T$ of the $d$ axis lengths to be $r/2$ and the others to
be $r$.  This gives
\[
 t_i\in\{r/2,r\},\qquad \prod_i t_i=T,\qquad t_d=r.
\]
Indeed $d\ell-\log_2T$ is an integer in $[0,d)$, so at least one axis
has length $r$; choose it as the last axis.
These definitions give
\begin{equation}\label{eq:sizes}
 Tp=\Theta(n),\quad d=\Theta(p^\epsilon),\quad
 K=\Theta(p^{\epsilon c}),\quad
 \ell=\Theta(p^{1-\epsilon}),\quad r=2^{\Theta(p^{1-\epsilon})}.
\end{equation}
The comparison constants can be fixed for every input length at once.
For sufficiently large $b$, the bounds on $n,T$ give
$b/2\leq\log_2T\leq b$. Since $p=6b$, $d=\lfloor b^\epsilon\rfloor$
and $d\leq b$, they imply
\[
 \frac1{12}p^\epsilon\leq d\leq p^\epsilon,
 \qquad \frac{p}{12d}\leq\ell\leq\frac{p}{3d}.
\]
Thus we fix $a_d=a_r=1/12$, $b_d=1$, $b_r=1/3$, and $C_\ell=1$
in the layer and transform interfaces. Their pointwise bounds cover every
computed pair $(d,r)$, including all choices arising at the same $p$.
These bounds also give that $r$ exceeds every fixed power of $p$,
$K/\log p\to\infty$, and
$\ell/K\to\infty$.  There are eventually enough axes for the
$O(\log d)$ reserved row axes and enough bits for all work blocks.

Set
\[
 \alpha=\left\lceil(12d^2b)^{1/4}\right\rceil,
 \qquad \gamma=2d\alpha^2,\qquad \eta=\frac1{4d}.
\]
The ceiling does not affect the needed bounds: for $A\ge1$,
$\lceil A\rceil\le2A$, and hence
\begin{equation}\label{eq:gamma}
 \alpha\le2(12d^2b)^{1/4},\qquad
 \gamma<28d^2\sqrt b\le28b^{1/2+2\epsilon}\le28b^{2/3}.
\end{equation}
It follows that $2\le\alpha<\sqrt p$ eventually and $\gamma=o(b)$.
For example $b\ge2^{24}$ implies $\gamma\le b/4$.

\subsubsection{Prime source lengths}

The distinct prime lengths $s_i$ must be close to the powers of two
$t_i$: their product must occupy a fixed fraction of the padded box,
while each ratio $t_i/s_i$ must leave enough separation for resampling.
We verify both requirements for the parameters just chosen.

The prime-interval result
\cite[Lemma~5.1]{HarveyHoeven2021}, based on explicit estimates of
Rosser and Schoenfeld~\cite{RosserSchoenfeld1962}, gives the following
bound: for
$0<\eta<1/4$ and $x>e^{2/\eta}$, the interval
\[
 (1-2\eta)x<q\leq(1-\eta)x
\]
contains at least $\eta x/(2\log x)$ primes. Here $\log$ is the natural
logarithm. The size bounds in \eqref{eq:sizes} give
$d=\Theta(p^\epsilon)$ and $\log r=\Theta(p/d)$.
For large $n$ we have $d\geq2$, so
$0<\eta<1/4$, and
\[
 \frac{\log(r/2)}d=\Theta(p^{1-2\epsilon})\longrightarrow\infty.
\]
Thus both $x=r/2$ and $x=r$ satisfy $\log x>8d=2/\eta$.
For either choice, the stated lower bound for the number of primes is
$x/(8d\log x)$ and exceeds $d$:
indeed, $\log x=\Theta(p^{1-\epsilon})$, whereas
$\log(8d^2\log x)=O(\log p)=o(\log x)$, so $x>8d^2\log x$ eventually.
The two intervals are disjoint because
$(1-\eta)r/2<(1-2\eta)r$ for $\eta<1/4$, and both intervals lie above
$2$. We may therefore choose distinct odd primes satisfying
\[
 (1-2\eta)t_i<s_i\leq(1-\eta)t_i.
\]
Every $\gcd(s_i,t_i)$ equals one because $t_i$ is a power of two.
Put $S=\prod_i s_i$. Multiplying the lower bounds on $s_i$ and
applying strict Bernoulli inequality gives
\[
 \frac ST>
 \left(1-\frac1{2d}\right)^d>\frac12
 \qquad(d\geq2).
\]
Together with $s_i<t_i$, this gives $T/2<S<T$.
Writing $\theta_i=t_i/s_i-1$, the upper bound on $s_i$ gives
$\theta_i\geq\eta/(1-\eta)>\eta$. Therefore
\[
 \alpha^4\theta_i>\frac{12d^2b}{4d}=3db\geq6b=p.
\]
The choice of $\alpha$ gives
$\alpha=\Theta(p^{1/4+\epsilon/2})$, and \eqref{eq:gamma} already
ensures $2\leq\alpha<\sqrt p$ for large $n$. We also have
$t_i\leq r\leq T<8n/b<2^{6b}=2^p$.
Thus every hypothesis of Lemma~\ref{lem:no-sort-resampling} holds for
all sufficiently large input lengths. The superpolynomial growth of $r$
established above also supplies the setup condition in
Lemma~\ref{lem:tensor-resampling}.

The machine finds these primes by scanning the two intervals
deterministically and testing each candidate by trial division through
its square root. Testing every candidate would visit $O(r)$ candidates
and at most $O(\sqrt r)$ trial divisors per candidate. Elementary
division on $O(\log r)$-bit integers and management of a list of at most
$d$ selected primes give the conservative bound
$O(dr^{3/2}\operatorname{poly}(p))$. Its logarithm is
\[
 O(\log p+\log r)=O(\log p+p/d)=o(p).
\]
Since $p=6\lceil\log_2 n\rceil=\Theta(\log n)$, the search takes
$n^{o(1)}=o(n)$ time.
All remaining scalar setup, including comparisons for fixed rational
powers, modular inverses,
prime products and shape descriptors, has cost polynomial in $p$.
For example, $d=\lfloor b^{2^{-75}}\rfloor$ is found by binary search
using exact comparisons $d^{2^{75}}\le b$; the exponent is fixed.
The same procedure computes $K$ and all other fixed rational powers.
All primes and length parameters are therefore generated by the machine;
they are not advice.

\subsubsection{The radix-digit polynomials}

Write the two inputs as polynomials in $2^b$, each with
$q=\lceil n/b\rceil$ nonnegative digits.  Their product has degree at most
\[
 2q-2<2n/b\le T/2<S.
\]
Thus convolution modulo $x^S-1$ has no nonzero wraparound.  Every product
coefficient obeys
\begin{equation}\label{eq:capacity}
 0\le c_j\le q(2^b-1)^2<2^{3b},
\end{equation}
since $q\le n\le2^b$.

\subsection{An explicit cyclic-convolution layout}

To use the source tensor, we first place the ordinary coefficient stream
in its tensor order. Agarwal and Cooley used the Chinese remainder
theorem to convert one cyclic convolution into multidimensional cyclic
convolution~\cite{AgarwalCooley1977}. The following map uses that
algebraic split, with coordinate powers that make the change of order a
sequence of controlled cyclic shifts.

Let $s_1,\ldots,s_d$ be the pairwise coprime
source lengths, let $S=\prod_i s_i$, and put
\[
 P_i=\prod_{j<i}s_j,\qquad \mu_i=P_i^{-1}\pmod{s_i}.
\]
Here $P_i$ denotes an integer prefix product, not a Fourier-coordinate
permutation.  The assignment
\[
 \Phi:\mathbb C[x]/(x^S-1)\longrightarrow
 \mathbb C[x_1,\ldots,x_d]/(x_1^{s_1}-1,\ldots,x_d^{s_d}-1),
 \qquad x\longmapsto\prod_i x_i^{\mu_i}
\]
respects $x^S=1$, so it defines an algebra map.  The target monomial
group has size $S$, and the image of $x$ has order $S$, because each
$\mu_i$ is a unit and the $s_i$ are pairwise coprime.  Its powers
therefore run through all target basis monomials exactly once.  Hence
$\Phi$ is an isomorphism and preserves cyclic convolution.

Every ordinary coefficient index has a unique mixed-radix expression
\[
 k=\sum_{j=1}^d a_jP_j,\qquad 0\le a_j<s_j.
\]
The coordinate of $\Phi(x^k)$ in axis $i$ is
\begin{equation}\label{eq:triangular-crt}
 b_i=a_i+\mu_i\sum_{j<i}a_jP_j\pmod{s_i}.
\end{equation}
Indeed terms with $j>i$ are divisible by $s_i$, and $\mu_iP_i=1$ modulo
$s_i$.  Numerical order of $k$ is lexicographic order of
$(a_d,\ldots,a_1)$.  Insert the padding in one nested scan of the full
box $\prod_i[0,t_i)$ in this lexicographic order, with $a_d$ the
outermost counter and $a_1$ the innermost.  At an address satisfying
$a_i<s_i$ for every $i$, copy the next original coefficient record;
at every other address, write a zero record.  The valid addresses
occur in the original coefficient order, so each record is copied to
its intended mixed-radix address.  To delete the padding after the
inverse axis reorder, use the same enumeration, copying the record at
each valid address and skipping the record at each invalid address.

Store the nested binary counters consecutively on a fixed work tape.
Their total width is $\sum_i\log_2t_i=\log_2T=O(p)$, and the stored
bounds $s_i$ also have total width $O(p)$.  At each of the $T$
addresses, updating the counters, testing validity, and handling the
$O(p)$-bit record therefore cost $O(p)$ steps, including counter-head
returns.  Each insertion or deletion scan costs $O(Tp)$.

Now reorder the complete padded axes to $(a_1,\ldots,a_d)$.  In this
physical order, the coordinates with
indices below $i$ precede coordinate $i$.  Update the targets in the
order $i=d,d-1,\ldots,1$ using \eqref{eq:triangular-crt}.  When target
$i$ is updated, none of its lower-index controls has yet changed, so
the offset is computed from the original digits $a_j$ required by the
formula.

To specify the action on padded addresses, let
\[
 \mathcal P=\prod_i[0,t_i),\qquad \mathcal V=\prod_i[0,s_i)
\]
be sets of integer addresses.  On a fixed control prefix, rotate the
valid target interval $[0,s_i)$ by the offset in
\eqref{eq:triangular-crt} and fix $[s_i,t_i)$.  If any lower-index
control is invalid, use the identity.  This defines a bijection of
$\mathcal P$: for fixed controls and spectators it is a cyclic
permutation of the valid target interval and the identity elsewhere.
Each step preserves $\mathcal V$ and its complement.  Consequently
zero padding remains zero after every step.  An inverse is explicit:
process $i=1,\ldots,d$, subtracting the same expression formed from the
already recovered lower coordinates; then undo the axis reorder and
perform the deletion scan just described.

Each triangular update is a prefix-controlled cyclic shift.  The
controls precede its target in the current physical order.  For each
prefix, prepare the offset, split the valid target interval into its
two consecutive pieces, and merge those pieces in the opposite order.
Every target position carries its full spectator suffix as a block;
the invalid interval is copied unchanged.  There are at most $T/t_i$
control prefixes for target $i$.  For any fixed setup degree $C$, all
offsets therefore cost $O(dTp^C/r)=o(Tp)$ to prepare, since
$t_i\ge r/2$ and $r$ exceeds every fixed power of $p$.

The axis widths are $\ell$ or $\ell-1$.  To exchange two adjacent axes
of unequal width, move the longer axis's leading bit outside their
equal-width pair, swap that pair by the proved chunk-swap procedure,
and put the leading bit at its new position.  This adds a fixed number
of linear passes to the equal-width swap cost.  Reversing all axes
uses $O(d^2)$ such moves.  Therefore the map $\Phi$, its inverse, and
their padding operations cost
\[
 O\bigl(Tp\,d^2(1+\ell^\tau)\bigr).
\]
This construction describes payload movement from address $a$ to its
new address.  It is distinct from the operator convention
$(R_cz)_j=z_{cj}$ used for Fourier dilations, which moves a payload
from address $j$ to $c^{-1}j$.

\subsection{Complex convolution by a coefficient twist}

The power-of-two complex convolution used in the chirp identity is
obtained from the synthetic convolution by a coefficient twist.
Put $\mathcal R=\mathcal R_r=\mathbb C[y]/(y^r+1)$ and
$\zeta=e^{\pi i/r}$.  The map
\[
 \Theta:\mathbb C[x_d]/(x_d^r-1)\longrightarrow\mathcal R,
 \qquad x_d\longmapsto\zeta y
\]
is well defined, since $(\zeta y)^r=(-1)(-1)=1$.  On the bases
$(x_d^k)_{0\le k<r}$ and $(y^k)_{0\le k<r}$ it is diagonal,
multiplying coefficient $k$ by the nonzero number $\zeta^k$.  It is therefore an
isomorphism, with inverse multiplying that coefficient by
$\zeta^{-k}$.  Both exact maps preserve the coefficient sup norm.
For an explicit check on the wrap, write $k+l=j+hr$, where
$h\in\{0,1\}$.  A product term after twisting and negacyclic reduction
has factor
\[
 \zeta^{k+l}(-1)^h=\zeta^j(\zeta^r)^h(-1)^h=\zeta^j.
\]
Untwisting therefore recovers the cyclic coefficient with the correct
sign.  We next specify how this linewise map meets the ring-convolution
interface.  Put $M=\prod_{i<d}t_i=T/r$.  For a complex array $f$ on
$\prod_{i=1}^d\mathbb Z/t_i\mathbb Z$, with $t_d=r$, define
\[
 (\iota f)_{\boldsymbol j}
   =\sum_{k=0}^{r-1}f_{\boldsymbol j,k}y^k,
 \qquad
 \boldsymbol j\in\prod_{i<d}\mathbb Z/t_i\mathbb Z.
\]
This identifies the scalar tensor with an array of $M$ records in
$\mathcal R$. As maps of values, $\iota$ and $\iota^{-1}$ preserve norms
and grids. Write
$(D_\pm f)_{\boldsymbol j,k}=\zeta^{\pm k}f_{\boldsymbol j,k}$.
The linewise extension of $\Theta$ is $\iota D_+$, with inverse
$D_-\iota^{-1}$.  Consequently, with $\mathcal M_r$ from
Section~\ref{sec:exact-ring-products}, normalized complex cyclic
convolution is
\[
 \mathcal M_{\mathbb C}(f,g)
   :=D_-\iota^{-1}
       \mathcal M_r(\iota D_+f,\iota D_+g)
    =\frac{f*g}{T}.
\]
Here $\mathcal M_r$ convolves the first $d-1$ axes over $\mathcal R$,
whereas $*$ on the right convolves all $d$ scalar axes.  The division
by $T$ is already part of $\mathcal M_r$.  The maps $D_\pm$ and
$\iota^{\pm1}$ are isometries, so $\mathcal M_{\mathbb C}$ is also a
bilinear contraction.

\subsubsection{Computed complex convolution}

For the tape representation, use the scalar format of
Section~\ref{sec:resampling}, with $w_{\rm c}=p+2$ bits per component.
For $p\geq2$, this satisfies $p+2\leq w_{\rm c}\leq2p$, so fix $C_w=2$
and instantiate the synthetic component width as $w=w_{\rm c}$.
Both representations then concatenate exactly the same real and
imaginary words in the same coefficient order: the last coordinate is
already the contiguous polynomial suffix. Hence $\iota$ and
$\iota^{-1}$ change only the shape descriptor and require no payload
movement when these component formats agree.

For the synthetic transform calls, the machine writes the header
\[
 \mathcal T=\Gamma(p)\Gamma(d)\Gamma(r)\Gamma(K)\Gamma(w_{\rm c})
       \prod_{i=1}^{d-1}\Gamma(t_i).
\]
For each normalized polynomial product it supplies
$\Gamma(p)\Gamma(r)\Gamma(w_{\rm c})$. These are precisely the headers
of Lemmas~\ref{lem:synthetic-transform-cost} and
\ref{lem:signed-ring-product}, formed from the computed parameters.
Their lengths are $O(p)$; forming the transform header costs
$\operatorname{poly}(p)=o(Tp)$, and copying a product header for each of
the $M=T/r$ polynomial records costs $O(Mp)=O(Tp)$.

For $0\le k<r$, obtain disk-grid approximations
$\widetilde\theta_k^+$ and $\widetilde\theta_k^-$ to $\zeta^k$ and
$\zeta^{-k}$, each with error less than $2\,2^{-p}$.  The exponential
routine uses denominator $2r$ and costs $O(p^{1+\delta})$ per factor.
Write $\widetilde{\mathcal M}_r$ for the procedure in
Proposition~\ref{prop:synthetic-convolution}.  For disk-grid inputs
$f,g$, one of norm at most $.7$, the computed complex convolution is the
following sequence,
with the twist arrays repeated on every last-coordinate line:
\[
 \begin{aligned}
 f'&=Q_p(\widetilde\theta^+\mathbin{\cdot}f),\\
 g'&=Q_p(\widetilde\theta^+\mathbin{\cdot}g),\\
 h'&=\widetilde{\mathcal M}_r(\iota f',\iota g'),\\
 \widetilde{\mathcal M}_{\mathbb C}(f,g)
   &=Q_p\!\left(\widetilde\theta^-\mathbin{\cdot}
                    \iota^{-1}h'\right).
 \end{aligned}
\]
Each phase product is computed exactly as a dyadic product before
$Q_p$ truncates its two components. The factors and disk operands have
numerators of magnitude at most $2^p$; four signed integer products and
their sums therefore use $O(p)$ bits. Truncation shifts the absolute
component numerators right by $p$ places and restores their signs, then
writes the resulting disk-grid value in the $w_{\rm c}$-bit component
format. It costs $O(p^{1+\delta})$ per coefficient. Generating every
factor anew therefore still gives total
twist cost $O(Tp^{1+\delta})$. The precision calculation below proves
that these calls return disk-grid arrays.

At the $h'$ step, Proposition~\ref{prop:synthetic-convolution} returns
its computed grid numerators after the exact scale $M$, in the known width
$w_M=w_{\rm c}+\log_2M$. The disk margin proved below gives
$\|h'\|<.701$ in every use here, so each component numerator has
magnitude at most $2^p$. Before applying $\iota^{-1}$, scan each
$w_M$-bit word and remove only leading sign-extension bits, writing the
same integer in $w_{\rm c}$ bits. The retained sign bit agrees with every
removed bit by that magnitude bound. This exact normalization costs
$O(Tw_M)=O(Tp)$ and supplies the matching scalar format for the output
untwist; it introduces no numerical truncation.

\subsection{Chirp and source-transform procedures}

We next use the complex convolution to evaluate the middle factor
$QF_{\boldsymbol t}$ in \eqref{eq:no-sort-tensor}, and then compose it
with the Gaussian maps. For the chirp array $a$ in
\eqref{eq:permuted-chirp}, the phase numerator
at tensor address $j$ is computed modulo $2r$ as
\[
 -\sum_i s_i j_i^2(r/t_i)\pmod{2r}.
\]
All integers have $O(p)$ bits and $r/t_i\in\{1,2\}$.  Scan the stored
axis descriptors and compute the $d$ squares and products with the
previously established integer multiplier.  Additions, reductions
modulo a power of two, and restarting local counters cost $O(p)$ per
term; multiplication costs $O(p\log p)=O(p^{1+\delta})$.  Exponential
evaluation then has the same $O(p^{1+\delta})$ bound.  Thus chirp
construction costs $O(Td p^{1+\delta})$ in total.  In particular, this
cost includes the scan of the axis descriptors at every coefficient.

Let $\widetilde a$ be the resulting disk-grid array, so
$\|\widetilde a-a\|_\infty<2\,2^{-p}$.  For a disk-grid target tensor
$x$ of norm at most $.6$, the computed chirp transform forms the arrays
\[
 \begin{aligned}
 x_0&=Q_p(\overline{\widetilde a}\mathbin{\cdot}x),\\
 y_0&=\widetilde{\mathcal M}_{\mathbb C}(\widetilde a,x_0),\\
 \widetilde{\mathscr C}(x)
     &=Q_p(\overline{\widetilde a}\mathbin{\cdot}y_0).
 \end{aligned}
\]
Thus $\widetilde a$ is the convolution kernel and $x_0$ is the data
operand after the first endpoint phase.  Replacing the approximations
by exact phases and exact convolution gives
$\overline a\mathbin{\cdot}\mathcal M_{\mathbb C}
(a,\overline a\mathbin{\cdot}x)=QF_{\boldsymbol t}x$ by
\eqref{eq:permuted-chirp}.  Below we show that $\|x\|_\infty\le.6$
implies $\|x_0\|_\infty\le.6$, which supplies the smaller operand
required by the computed convolution.

\subsubsection{The source transform}

For a source disk-grid tensor of norm at most $.5$, write
$\widetilde{\mathscr F}$ for the following procedure: apply the linewise
approximation to $\mathcal A$, apply the
computed chirp transform $\widetilde{\mathscr C}$, apply the linewise
approximation to $\mathcal B_0$, and multiply the grid numerators
exactly by $2^\gamma$.  Its exact counterpart is
\[
 \mathscr F=2^\gamma\mathcal B_0QF_{\boldsymbol t}\mathcal A
            =RF_{\boldsymbol s}.
\]
Here an exact counterpart replaces every numerical subroutine by the
exact map it approximates and omits truncations.  The source arrays are
stored in the padded power-of-two box, with zeros outside the valid
source coordinates.
The compressed scalar output uses $w_{\rm c}$-bit components. Appending
$\gamma$ zero bits to each word gives the exact scaled numerator in the
known signed width $w_{\rm c}+\gamma$. For the stated input norm, the
disk estimate below proves that every scaled coordinate still belongs to
$\mathbb D_p$. Its component numerators therefore have magnitude at most
$2^p$, so a scan removes only leading sign-extension bits and returns
the same integers in $w_{\rm c}$ bits. This scan costs
$O(T(w_{\rm c}+\gamma))=O(Tp)$ and leaves the padded zero records zero.
Define the opposite-sign procedure by
\[
 \widetilde{\mathscr F}^{+}(z)
   =\overline{\widetilde{\mathscr F}(\overline z)};
\]
its exact counterpart is $RF_{\boldsymbol s}^{+}$ by
Lemma~\ref{lem:chirp-permutation}.

\subsection{The multiplication algorithm}

The source-transform procedures now supply the three transforms needed
for the radix-digit convolution. Read each stored operand backward from
its least significant end in groups of $b$ bits. This produces the radix-$2^b$ digits
$a_0,\ldots,a_{q-1}$ in increasing coefficient order.  Buffer each
group on a fixed work tape and read it back most significant bit first,
reversing its least-significant-first input order; pad the final short
group with leading zeroes.
For digit $a_j$, write the exact grid numerator
$a_j2^{p-b-2}$ in a signed $w_{\rm c}$-bit real word, followed by a zero
$w_{\rm c}$-bit imaginary word. This exponent is nonnegative
because $p=6b$, and the represented value is exactly $a_j2^{-b-2}$.
Append $S-q$ zero coefficient records.  Buffering and ordering the
groups costs $O(n)$, including the one partial group's padding, and
writing the $S$ records costs $O(Sp)$.
Together with the initial count and copies, this input conversion
costs $O(n+Sp)=O(Tp)$; all $n$ input bits, including leading zeroes,
are retained before the final group is padded.

With these streams prepared, perform the following operations.
\begin{enumerate}
\item Apply the coefficient-address map $\Phi$, with its padded layout, to obtain
source arrays $u,v$.  Their coefficients are exact $p$-bit grid
values of norm less than $1/4$.
\item Compute $\widetilde{\mathscr F}(u)$ and
$\widetilde{\mathscr F}(v)$.
\item Multiply corresponding complex coefficients exactly and truncate
each real and imaginary part toward zero to $p$ fractional bits, writing
the result in the common $w_{\rm c}$-bit component format.
Denote the resulting source array by $z$.
\item Compute $\widetilde{\mathscr F}^{+}(z)$ and multiply its grid
numerators exactly by $S$, retaining the denominator $2^p$ and writing
each component in signed width $w_{\rm c}+b$. Denote the represented
dyadic array by $\widetilde w$.
\item Multiply the grid numerators of $\widetilde w$ exactly by
$2^{2b+4}S$, again retaining the denominator $2^p$, and write each
component in signed width $w_{\rm c}+4b+4$. Round the represented real
values to the nearest integers. At each valid address, write the rounded
coefficient in one signed $(3b+1)$-bit word, most significant bit first;
at every padded address,
write a zero word of that width. Supply this local payload width to
$\Phi^{-1}$ and apply it, including removal of the address padding, to
restore the ordinary coefficient order.
\item Propagate radix-$2^b$ carries in increasing coefficient order.
Write the resulting bits from least to most significant on a work tape,
then reverse that stream and output exactly $2n$ bits in
most-significant-first order.
\end{enumerate}

The two occurrences of $S$ in this algorithm undo different
normalizations.  By \eqref{eq:source-permutation-cancel}, the exact
counterpart after the three transforms is $(u*v)/S^2$, so after the
first scale it is $w=(u*v)/S$.  The second scale gives
$2^{2b+4}(u*v)$, whose inverse image under $\Phi$ is the original
integer coefficient convolution.  The next subsection verifies the
disk bounds needed by every call and proves that the final error is
less than $1/2$; it also justifies the rounding, carries, and output
length.

\subsection{Precision and exact recovery}

Proposition~\ref{prop:simultaneous-layer} uses
$C_1=20$, so $\epsilon C_1<1$ makes its coefficient width $O(p)$.
Lemma~\ref{lem:synthetic-transform-cost} therefore applies to every
synthetic transform here, with error less than
$\sqrt2\ell\,2^{-p}<L2^{-p}$, where $L=\log_2T$, for all
sufficiently large $n$.

\subsubsection{Scales and disk margins}

Put $h_p=2^{-p}$, and use the coefficient sup norm throughout this
subsection.  An error of $a$ will mean absolute error less than $ah_p$.
A linear contraction propagates an input error without increasing it.
For a bilinear contraction $\mathcal M$, replacing its operands gives
\[
 \|\mathcal M(f',g')-\mathcal M(f,g)\|
 \le \|f'-f\|\,\|g'\|+\|f\|\,\|g'-g\|.
\]
When all four operands are in the disk, their two errors therefore add.

We also need the norm as well as the error of a computed phase product.
Let $e$ be a unit phase, let $\widetilde e$ be a disk-grid approximation
with $|\widetilde e-e|<2h_p$, and let $x$ be a disk operand.  The
truncation bound in \eqref{eq:transform-truncate} gives
\[
 \|Q_p(\widetilde e x)-ex\|
   <(2\|x\|+\sqrt2)h_p<4h_p,
 \qquad
 \|Q_p(\widetilde e x)\|\le\|x\|.
\]
The second inequality follows from $|\widetilde e|\le1$ and the fact
that $Q_p$ does not increase the norm.  It preserves every smaller
operand norm through a twist or an endpoint phase multiplication.
The same bounds hold coordinatewise for arrays of phases.

For disk-grid ring arrays $f,g$, put $M=T/r$.  Two normalized synthetic
forwards, pointwise polynomial multiplication divided by $r$, and the
normalized opposite transform give $(f*g)/(rM^2)$.  The exact
multiplication by $M$ in Proposition~\ref{prop:synthetic-convolution}
then gives $\mathcal M_r(f,g)=(f*g)/T$.  Before that multiplication,
the three transform errors and the polynomial-product truncation total
less than $3L+2$.  Thus the error of the returned normalized
convolution is bounded by
\[
 E_{\rm ring}:=M(3L+2)<5TL.
\]
If one operand has norm at most $.7$, bilinear contraction gives
$\|\mathcal M_r(f,g)\|\le.7$.  Its computed grid output consequently
has norm less than $.7+5TLh_p<.701$ for large $n$.  This proves the
disk margin after the scale $M$, without clipping the output.

Consider now $\widetilde{\mathcal M}_{\mathbb C}(f,g)$ for disk grid
operands, one of norm at most $.7$.  Its two computed input twists
$f',g'$ are disk grid arrays, and the smaller one still has norm at
most $.7$ by the phase inequality.  The identification $\iota$
preserves these bounds.  The preceding ring margin therefore gives
$\|h'\|<.701$.  Neither $\iota^{-1}$ nor the computed output untwist
increases this norm.  In particular,
$\widetilde{\mathcal M}_{\mathbb C}(f,g)$ is a disk grid array.

To compare it with $\mathcal M_{\mathbb C}(f,g)$, insert the exact
normalized $\mathcal M_r$ at $\iota f'$ and $\iota g'$.
Its bilinear contraction bound propagates the two input twist errors,
each less than $4$.  Its approximation contributes $E_{\rm ring}$,
and the output untwist contributes less than $4$.  All these
comparisons are at the normalized $\mathcal M_r$ interface, after its
factor $M$.  They give
\[
 \|\widetilde{\mathcal M}_{\mathbb C}(f,g)
            -\mathcal M_{\mathbb C}(f,g)\|
 <(E_{\rm ring}+4+4+4)h_p<6TLh_p.
\]

For a disk-grid chirp input $x$ with $\|x\|\le.6$, the previously
defined data operand $x_0=Q_p(\overline{\widetilde a}\cdot x)$ has
norm at most $.6$.  The kernel $\widetilde a$ is a disk grid array.
Thus this pair satisfies the smaller-operand condition just proved:
$\|y_0\|<.701$, and the last endpoint multiplication gives
$\|\widetilde{\mathscr C}(x)\|<.701$.  In particular, the chirp output
is a legal disk-grid input to the compression maps.

For its error, the exact operands $a$ and $\overline a\cdot x$ are in
the disk because every entry of $a$ has modulus one.  The kernel
approximation contributes less than $2$ and
the first endpoint multiplication contributes less than $4$ when
inserted into the exact bilinear contraction
$\mathcal M_{\mathbb C}$.  The computed convolution contributes less
than $6TL$, and the last endpoint multiplication less than $4$.
Equation~\eqref{eq:permuted-chirp} identifies the exact result, so
\[
 \|\widetilde{\mathscr C}(x)-QF_{\boldsymbol t}x\|
 <(6TL+2+4+4)h_p<8TLh_p.
\]

Finally let the input to a source transform have norm at most $.5$.
The computed expansion has error less than $dp^2$ and norm less than
$.5+dp^2h_p<.51$.  It is therefore a permitted chirp input.  The
chirp output is in the disk by the preceding argument, so each
compression call is defined and returns another disk-grid array.
Insert the exact contractions $\mathcal A$, $QF_{\boldsymbol t}$,
and $\mathcal B_0$ at the computed inputs.  The error before the
source scale is less than $2dp^2+8TL$.  Multiplication by
$2^\gamma$ therefore gives error less than
$2^\gamma(2dp^2+8TL)h_p$, which is bounded by $E_sh_p$ with
\begin{equation}\label{eq:source-error}
 E_s=9\,2^\gamma TL,
\end{equation}
because $2dp^2<TL$ eventually.  The exact source transform
$RF_{\boldsymbol s}$ is a contraction, so the scaled approximation
has norm at most $.5+E_sh_p<.501$.  Its exact numerator scaling
therefore also leaves a disk-grid output without clipping.
Conjugation preserves the grid and norm, so the opposite procedure
has the same error and disk guarantee.
All these margins may be enforced by one fixed cutoff.  For example,
$b\ge2^{24}$, together with the other eventual construction conditions,
makes $dp^2h_p$, $8TLh_p$, $E_sh_p$ and
$28\,2^\gamma T^2Lh_p$ smaller than $10^{-3}$.

\subsubsection{The final integer coefficients}

The scaled radix digits in $u,v$ are exact $p$-bit grid values of norm
strictly less than $1/4$.  Put
\[
 \begin{aligned}
 U&=RF_{\boldsymbol s}u,&\qquad V_1&=RF_{\boldsymbol s}v,\\
 \widetilde U&=\widetilde{\mathscr F}(u),&
 \widetilde V_1&=\widetilde{\mathscr F}(v).
 \end{aligned}
\]
The exact arrays have norm less than $1/4$, and the computed arrays
have norm less than $1/4+E_sh_p<.26$.  Hence
$z=Q_p(\widetilde U\cdot\widetilde V_1)$ has norm at most
$.26^2<.5$.  To estimate its error, write
\[
 \widetilde U\cdot\widetilde V_1-U\cdot V_1
 =(\widetilde U-U)\cdot\widetilde V_1
       +U\cdot(\widetilde V_1-V_1).
\]
The factors $\widetilde V_1$ and $U$ are in the disk.  The two terms
have total norm less than $2E_sh_p$, and truncation adds less than $2h_p$.
Thus $\|z-U\cdot V_1\|<(2E_s+2)h_p$.

Let $G:=RF_{\boldsymbol s}^+$ be the exact opposite contraction.
Lemma~\ref{lem:chirp-permutation}, applied to
$U=RF_{\boldsymbol s}u$ and $V_1=RF_{\boldsymbol s}v$, gives
\[
 G(U\cdot V_1)=(u*v)/S^2.
\]
The computed opposite transform may be applied to $z$ because
$\|z\|<.5$. Insert $G$ at that computed input:
\[
 \begin{aligned}
 \|\widetilde{\mathscr F}^{+}(z)-G(U\cdot V_1)\|
 &\le \|\widetilde{\mathscr F}^{+}(z)-Gz\|\\
 &\qquad+\|G(z-U\cdot V_1)\|\\
 &<(3E_s+2)h_p.
 \end{aligned}
\]
The first exact numerator multiplication by $S$ gives
$\widetilde w$, approximating $w=(u*v)/S$ with error less than
$E_wh_p$, where
\[
 E_w:=S(3E_s+2)<28\,2^\gamma T^2L.
\]
Each coefficient of $w$ is an average of $S$ products of coefficients
of $u$ and $v$, so $\|w\|<1/16$.  The earlier cutoff therefore gives
\[
 \|\widetilde w\|<1/16+E_wh_p<1/16+10^{-3}<1.
\]
This proves the disk bound for the computed first scaled result.

In tensor order, the unscaled integer coefficient array is
$2^{2b+4}S w$.  This second occurrence of $S$ reverses the
normalization still present in $w$, and $2^{2b+4}$ reverses the two
input scales.  Both multiplications act exactly on grid numerators.
Since $S\le T<2^b$ and $L<b$ eventually, the final absolute error is less than
\begin{equation}\label{eq:final-error}
 2^{2b+4}S\cdot28\,2^\gamma T^2L\cdot2^{-6b}
 <448b\,2^{-b+\gamma}
 \le448b\,2^{-3b/4}<\frac12.
\end{equation}
The true coefficients are real integers.  If $H_j$ is the real
numerator after the second scaling, its represented value is
$H_j/2^p$.  The strict error below $1/2$ makes the nearest integer to
that value the unique correct coefficient.

These exact operations use known fixed record widths. If $N$ is a
signed $w_0$-bit integer, then $-2^{w_0-1}\leq N<2^{w_0-1}$. Since
$0<S<2^b$, its product satisfies
$-2^{w_0+b-1}<SN<2^{w_0+b-1}$ and therefore fits in signed width $w_0+b$.
The first product by $S$ starts at width $w_{\rm c}$ and is written in
width $w_{\rm c}+b$. The next product by $S$ is written in width
$w_{\rm c}+2b$; appending $2b+4$ least significant zero bits then gives
the exact second scale in width $w_{\rm c}+4b+4$. Each complex record
still places its most-significant-bit-first real word before its imaginary
word. The known widths determine every boundary in the next scan and are
$O(p)$ because $p=6b$.

The established integer multiplier performs the two products by $S$ in
$O(p\log p)$ time per entry. To serialize each exact signed product at
its stated width, zero-pad its magnitude and complement and increment
for a negative sign. The proved bounds ensure that no significant bit is
discarded. These conversions and the power-of-two shift use linear scans.
To round $H_j/2^p$, recover the real magnitude, inspect its low $p$ bits,
and apply the rounding carry with the sign restored. The strict error
bound makes the result the correct nonnegative coefficient, and
\eqref{eq:capacity} puts it in $[0,2^{3b})$. It therefore fits in one
signed $(3b+1)$-bit two's-complement word, written most significant bit
first. Scan the padded box in the
physical axis order recorded by its shape descriptor, using the same
$O(p)$-bit nested counters and validity tests as above. Consume the
imaginary word and write this real coefficient word at each valid address,
with an all-zero
word at each padded address. This serialization costs $O(Tp)$ in total.
It keeps the tensor address order and supplies payload width $3b+1$ to
the local descriptors for $\Phi^{-1}$, so the address routine receives a
stream of known fixed-length records.

Lemma~\ref{lem:signed-ring-product} proves exact signed packing,
centered extraction and negacyclic reduction. Its only numerical error
is the stated truncation after division by $r$, already included above.
Undo the coefficient-address isomorphism $\Phi$ before propagating
radix carries.

To propagate carries, let $Q_0=2^b$.  The recovered ordinary
coefficients are the $c_j$ in \eqref{eq:capacity}; they are zero for
$2q-2<j<S$.  With $q=\lceil n/b\rceil$ as
above, use the recurrence for $0\le j<S$:
\[
 k_0=0,\qquad d_j=(c_j+k_j)\bmod Q_0,\qquad
 k_{j+1}=\left\lfloor\frac{c_j+k_j}{Q_0}\right\rfloor.
\]
The carries satisfy $0\le k_j\le q(Q_0-1)<Q_0^2$.  If this holds at
$j$, then
\[
 c_j+k_j\le q(Q_0-1)^2+q(Q_0-1)=qQ_0(Q_0-1),
\]
which proves the next carry bound.  The defining relation
$c_j+k_j=d_j+Q_0k_{j+1}$ also gives the exact invariant
\[
 \sum_{h<j}c_hQ_0^h
   =\sum_{h<j}d_hQ_0^h+k_jQ_0^j.
\]
It holds for $j=0$, and the defining relation proves its next instance.
The unprocessed terms are multiples of $Q_0^j$, so the identity fixes
the first $j$ radix digits of the full product.  The value $k_j$ is
the carry from the processed coefficients into the remaining positions.

After $\Phi^{-1}$ the coefficients occur in increasing ordinary index,
so this carry scan proceeds from least to most significant.  Carry
storage and work per coefficient are $O(b)$ bits and steps.  Write
each digit's bits from low to high on a work tape, and drain the final
$O(b)$-bit carry in the same order.  Copy the completed stream backward
once to obtain most-significant-first order.  The stream has length
$O(Sb)=O(Tp)$, including chunk padding and the final carry, so this
reversal is a charged linear tape pass.  The exact product is less
than $2^{2n}$; all bits above position $2n-1$ are therefore zero.
Remove only those zero padding bits and retain leading zeros as needed
to produce exactly $2n$ bits.  Zero inputs obey the same invariant and
output rule.

\subsection{The complete time bound}

For the cost calculation put $V=Tp=\Theta(n)$. The fixed-width coefficient
formats use $\Theta(p)$ bits per coefficient, so a full padded array has
$\Theta(V)$ stored bits. Table~\ref{tab:costs} lists costs after division
by $V$; a fixed number of repeated transforms and convolutions is
included in its constants.
\begin{table}[ht]
\centering
\small
\begin{tabular}{@{}>{\raggedright\arraybackslash}p{.40\linewidth}ll@{}}
\toprule
Operation & Normalized cost & Power of $p$\\
\midrule
Prefix-slot moves and individual butterfly rounds & $dK$ & $\epsilon(1+c)$\\
Chunk exchanges for transform layout & $pK^{\tau-1}$ & $1-\epsilon c(1-\tau)$\\
Simultaneous butterfly rounds & $\ell d^{\lambda'}$ & $1-\epsilon(1-\lambda')$\\
CRT and axis layouts & $d^2(1+\ell^\tau)$ & $\tau+\epsilon(2-\tau)$\\
Gaussian line maps & $dp^{1/2+\delta}\alpha$ & $3/4+\delta+3\epsilon/2$\\
Chirps, twists and scalar products & $dp^\delta$ & $\epsilon+\delta$\\
Packed polynomial products & $\log(rp)$ & $1-\epsilon$\\
Final scaling and rounding & $p^\delta$ & absorbed above\\
Input, padding, carries and output & $1$ & absorbed above\\
\bottomrule
\end{tabular}
\caption{Costs divided by $Tp$; each coefficient
representation has $O(p)$ bits.}\label{tab:costs}
\end{table}

The seven dominant powers in Table~\ref{tab:costs} have the following
margins below one:
\begin{equation}\label{eq:margin-list}
\begin{split}
 g_1&=1-\epsilon(1+c),\qquad
 g_2=\epsilon c(1-\tau),\qquad
 g_3=\epsilon(1-\lambda'),\\
 g_4&=1-\tau-\epsilon(2-\tau),\qquad
 g_5=1/4-\delta-3\epsilon/2,\\
 g_6&=1-\delta-\epsilon,\qquad g_7=\epsilon.
\end{split}
\end{equation}
Direct substitution gives
\[
 g_1>\frac12,\quad g_2=2^{-181},\quad g_3=2^{-129},\quad
 g_4>2^{-51},\quad g_5>\frac18,\quad g_6>\frac12,\quad g_7=2^{-75}.
\]
Thus $\min_i g_i=2^{-181}=2\kappa$.  This explicit slack will cover
any remaining fixed power of $\log p$.

Here $d^2$ is absorbed by $d^2\ell^\tau$, and
$\log(rp)=O(p^{1-\epsilon}+\log p)=O(p^{1-\epsilon})$.
The first three rows are the terms in
\eqref{eq:synthetic-transform-cost}. The prefix-slot moves and
individual butterfly rounds in Lemma~\ref{lem:synthetic-transform-cost}
use $O(dK)$ linear scans. Its positional layout uses $O(p/K)$ exchanges
of width-$K$ chunks, because $d\ell=O(p)$. Each exchange costs
$O(VK^\tau)$ by Lemma~\ref{lem:chunk-swap}, so the normalized total is
\[
 O(p/K)\,K^\tau=O(pK^{\tau-1}).
\]
For the at most $\ell$ simultaneous rounds,
Proposition~\ref{prop:simultaneous-layer} supplies $O(Vd^{\lambda'})$
per round. Its bound already includes the binary basis changes,
deterministic repair, row padding and removal, all base-$m$ groups,
and fixed-tape cleanup.

Axis exposure and restoration require $O(d^2)$ adjacent axis moves.
Each costs $O(V(1+\ell^\tau))$: peel a single excess bit when widths
differ by one, exchange their equal-width parts, and restore the bit.
The $d$ CRT rotations copy valid and invalid intervals in $O(dV)$ time.
Their offset setup was bounded above by $O(dTp^C/r)=o(Tp)$,
using at most $T/t_i$ control prefixes for target $i$.

For Gaussian resampling, axis $i$ has at most $T/t_i$ processed lines.
The line-machine costs therefore sum to
\[
 \sum_{i=1}^d\frac{T}{t_i}
       O(t_i p^{3/2+\delta}\alpha)
 =O(dTp^{3/2+\delta}\alpha).
\]
Dividing by $V=Tp$ gives the fifth row.  The factor
$\alpha=\Theta(p^{1/4+\epsilon/2})$ explains its displayed power of $p$.

Chirp numerators have $O(p)$ bits and can be computed modulo $2r$ as
sums of $d$ terms $s_i j_i^2(r/t_i)$.  Computing their squares and
products with the previously established $O(N\log(2N))$ multiplier
on $N$-bit inputs,
then generating the exponentials, costs $O(dp^{1+\delta})$ per entry.
The smaller twist and scalar pointwise-product costs fit the same row:
they use a fixed number of exact products on $O(p)$-bit component
numerators, followed by linear truncation scans.  The two final products
by $S$ also cost $O(p\log p)$ per entry, so their normalized cost is
$O(\log p)=O(p^\delta)$ as listed separately.  There are $T/r$
polynomial records in each pointwise ring product.  Four integer
products on $O(rp)$ bits per record therefore cost
$O(Tp\log(rp))$ in total.  Signed packing, centered extraction,
negacyclic wrapping, and copying are linear stream passes.  These calls
use the previously established multiplier on records of length
$O(rp)=o(n)$; they do not call the improved algorithm being constructed.
The exact source scaling by $2^\gamma$ is a shift of $O(p)$-bit
numerators by $\gamma=O(p)$ places, so it is a linear scan per entry
and fits the final-scaling row. The scalar-format adapters and the exact
removal of sign-extension bits after the ring and source scales also use
linear scans of their $O(p)$-bit components, within the listed scan costs.

\paragraph{Setup and exceptional streams across all calls.}
We next separate setup at recursion nodes from work repeated for every
polynomial record or every Gaussian line.  Here a visit means a node
of the layer or address-layout recursions, or one of their round and
group invocations.  The internal recursion of the established integer
multiplier and the iterations of the quoted Gaussian algorithms are
already included in their separately charged running times.
The layer and address-layout recursions have constant branching and
depth $O(\log p)$, and the number of their rounds and base-$m$ groups
is polynomial in $p$.  Choose a fixed exponent $A$ so that their
total number of visits is at most $p^A$, and a fixed $C$ bounding
descriptor setup performed once per visit by $p^C$ steps.  This setup then
costs at most $p^{A+C}=o(Tp)$.

At a layer or synthetic-transform visit that processes polynomial
records, there are at most $O(T/r)$ such records, including the
bounded row padding.  Every relevant child keeps the whole
length-$r$ polynomial suffix.  Consequently $O(p^C)$ descriptor or
address work per record, over these visits, costs at most
$O(Tp^{A+C}/r)=o(Tp)$.  Setup before each Gaussian line call has a
different count:
there are at most $2dT/r$ lines per tensor map and only a fixed number
of such maps, so $O(p^C)$ setup per line costs
$O(dTp^C/r)=o(Tp)$.  These estimates all use that $r$ exceeds every
fixed power of $p$.

The general address-swap subroutine also covers one-bit payloads,
where no polynomial suffix is available. For root volume $V_0$ and
current child volume $V_j$, the descriptor estimate in the proof of
Proposition~\ref{prop:power-interchange} gives
$(\log(2V_0))^{O(1)}=O(V_j)$. It uses the invariant that both original
chunk ranges remain present in every child, and therefore includes
polynomial descriptor work even for a one-bit payload.
Elementary counter updates use local length copies, so they do not
rescan a growing descriptor for each data bit.  The chirp computation
does scan the axis descriptors at every coefficient; its
$O(Tdp^{1+\delta})$ cost was included in the chirp row.

For one packed change of basis, let $V_{\rm call}$ be its current
logical volume.  The bad-address fraction is
$O(d2^{-K})$, independently of the payload.  Each polynomial record
has $\Theta(rp)$ payload bits, and an exceptional record adds an
$O(p)$-bit destination key.  Thus the compact exceptional stream has
length
\[
 O\!\left(\frac{V_{\rm call}}{rp}\,d2^{-K}(rp+p)\right)
   =O(V_{\rm call}d2^{-K}).
\]
Stable bit sorting makes $O(p)$ scans of this stream, for one-call
cost $O(V_{\rm call}dp2^{-K})$.  Extraction and reinsertion use a
constant number of scans of the full current volume; those scans are
already in the node overhead.  Since there are at most $p^A$ visits,
each of volume $O(V)$, the compact sorting costs sum to at most
\[
 O\!\left(dp2^{-K}\sum_{\rm calls}V_{\rm call}\right)
   =O(Vp^{A'}2^{-K})=o(V)
\]
for another fixed exponent $A'$.  Here $d=O(p^\epsilon)$ and
$K=\Theta(p^{\epsilon c})$ with $\epsilon c>0$.  Predicate and key
arithmetic is the polynomial work per record counted above.  This is
a deterministic worst-case bound.

Finally, all role, arithmetic, and stack tapes form fixed finite
families.  A child parks and restores only its parent's streams, with
the stack head at the active top; copying and cleanup are charged to
that logical volume.  No operation repeatedly scans parked ancestors.
Growing axis lists and call states are stored data, not extra tapes or
alphabet symbols.  Global setup and prime search cost $o(V)$, as
proved earlier.

Each of the seven displayed powers is at most $1-2\kappa$ by
\eqref{eq:margin-list}.  Any fixed power of $\log p$ is at most $p^\kappa$
eventually.  The complete large-input bound is consequently
\[
 O(Tp\,p^{1-\kappa})=O(n(\log n)^{1-\kappa}).
\]
Fix one cutoff beyond all the size, precision, margin and layout
conditions, and use schoolbook multiplication below it.  The finitely
many smaller lengths are included by enlarging the implicit constant.
With $\lg n=\max\{\lceil\log_2n\rceil,1\}$, the resulting single
fixed-alphabet, fixed-tape deterministic machine has worst-case time
$O(n(\lg n)^{1-\kappa})$ for every input length and produces precisely
the required $2n$-bit product.
